\section{Preliminary Analysis}



\noindent
\textbf{Notation}.
We follow the notation established in~\citep{kalai2025consensus}. For a prompt $x$, each model $i \in \{1, \dots, n\}$ induces a probability distribution $p_i(y)$ over a output space $\mathcal{Y}$. Let $\mathrm{Distr}(\mathcal{Y})$ denote the set of probability distributions on $\mathcal{Y}$. $\mathcal{Y}$ is the union of the safe space $S$ and the unsafe space $U$. For any $p \in \mathrm{Distr}(\mathcal{Y})$ and any subset $H \subseteq \mathcal{Y}$, we define the cumulative probability $p(H) = \sum_{y \in H} p(y)$. Given a collection of distributions $(p_1, \ldots, p_n) \in \mathrm{Distr}(\mathcal{Y})^n$, let $p_{(i)}(y)$ represent the $i$-th smallest of the probabilities $\{ p_t(y) \}_{t=1}^n$. We use the symbol $\bot \notin \mathcal{Y}$ to denote an abstention. 

\noindent
\textbf{Output Distribution}. Because the model's output is terminated either by special tokens or by a maximum token limit, for a finite tokenizer, We follow existing work \citep{chijiwa2025lossless} to treat the ostensibly unbounded output distribution as a finite set to facilitate analysis. Formally, for model $i$ with maximum token length $L$ and tokenizer $T_i$, the output space is $\mathcal{Y}_i = \bigcup_{j=1}^{L} T_i^j$.


\noindent
\textbf{Model Group}.
The model group ($\mathcal{MG}$) consists of $n$ generative models, including $s$ safe models and 
$f=n-s$ unsafe models. Each safe model should maintain a reasonable probability distribution for any input prompt $x$. The probability that the safety model outputs a safe response is $\Psi > 0$. Unsafe models are assumed to be fully controllable by an adversary and may exhibit arbitrary unsafe behaviors or follow any probability distribution, and we refer to this as the Byzantine model based on the definition of Byzantine replicas in classical BFT consensus research \citep{zhang2023waterbear, dashing2024duan, Das2024AC}.



